\section{Conclusions}
\label{sec:conclusions}

We introduced \textbf{DeltaKV}, a residual-based KV cache compression framework that exploits long-range inter-token similarity to remove redundant shared structure and retain only lightweight residuals. This design reduces KV memory usage to 29\% of the original size while preserving the information most relevant for attention. When co-designed with our \textbf{Sparse-vLLM} inference engine, DeltaKV delivers up to $2\times$ higher decoding throughput and near-lossless accuracy on diverse long-context and reasoning benchmarks, including LongBench and AIME.
%%%
Beyond empirical gains, we highlight two insights: (i) KV caches exhibit substantial global redundancy beyond local similarity, and (ii) residuals after reference subtraction have favorable statistics for compression and quantization. 
This suggests KV optimization should move beyond eviction and low-rank projection toward similarity-aware representations.

% Overall, DeltaKV demonstrates that a simple, GPU-friendly residual formulation, combined with system-level support for sparse and irregular KV layouts, can effectively alleviate memory and bandwidth bottlenecks in long-context LLMs. We hope this work motivates further research on similarity-aware caching, hybrid compression-sparsity methods, and co-design of algorithms and inference systems for scalable LLM deployment.



\section*{Impact Statement}
This work contributes to the efficiency of long-context LLM serving by significantly reducing memory footprints and increasing inference throughput. These advancements promote energy-efficient computing and broaden the accessibility of powerful models on resource-constrained hardware.

While increasing the efficiency of foundational models is generally beneficial, it may also lower the barrier for indiscriminate large-scale deployment, potentially amplifying existing societal risks associated with LLMs. Additionally, although our compression method demonstrates robust performance, reliance on approximated representations could theoretically impact model reliability in unforeseen edge cases. We encourage practitioners to conduct thorough evaluations and adhere to responsible deployment guidelines when applying these techniques in real-world systems.

% Authors are \textbf{required} to include a statement of the potential broader
% impact of their work, including its ethical aspects and future societal
% consequences. This statement should be in an unnumbered section at the end of
% the paper (co-located with Acknowledgements -- the two may appear in either
% order, but both must be before References), and does not count toward the paper
% page limit. In many cases, where the ethical impacts and expected societal
% implications are those that are well established when advancing the field of
% Machine Learning, substantial discussion is not required, and a simple
% statement such as the following will suffice:

% ``This paper presents work whose goal is to advance the field of Machine
% Learning. There are many potential societal consequences of our work, none
% which we feel must be specifically highlighted here.''

% The above statement can be used verbatim in such cases, but we encourage
% authors to think about whether there is content which does warrant further
% discussion, as this statement will be apparent if the paper is later flagged
% for ethics review.




% \section*{Accessibility}

% Authors are kindly asked to make their submissions as accessible as possible
% for everyone including people with disabilities and sensory or neurological
% differences. Tips of how to achieve this and what to pay attention to will be
% provided on the conference website \url{http://icml.cc/}.




% \section*{Software and Data}

% If a paper is accepted, we strongly encourage the publication of software and
% data with the camera-ready version of the paper whenever appropriate. This can
% be done by including a URL in the camera-ready copy. However, \textbf{do not}
% include URLs that reveal your institution or identity in your submission for
% review. Instead, provide an anonymous URL or upload the material as
% ``Supplementary Material'' into the OpenReview reviewing system. Note that
% reviewers are not required to look at this material when writing their review.




% % Acknowledgements should only appear in the accepted version.
% 